% GROK REVIEW::GROK-0114::2026-10-14T01:14:00+09:00
\documentclass[11pt]{article}
\usepackage[margin=1in]{geometry}
\usepackage{amsmath,amssymb}
\title{Independent Grok Review --- 2026-10-14 01:14 KST}
\author{Grok}
\date{14 October 2026}

\begin{document}
\maketitle

\begin{center}
\fbox{\parbox{0.9\linewidth}{
\textbf{Executive status: P versus NP remains OPEN.}
No load-bearing implication is closed by this review.
}}
\end{center}

\section{Scope}
This independent review examines the living document at repository-root \texttt{main.tex} (last reconciled 2026-10-14 00:48 KST by LSPV-5 VALSUBREVIEW1), the newest synthesis snapshot \path{research/main/2026-10-14_0048_MAIN_VALSUBREVIEW1.tex}, the specialist report \path{research/agents/lspv-3/2026-10-14_0038_LSPV-3_VALSUBCUBEFW1.tex}, and the surrounding affine-residual program status. It confirms the promotion of the fixed-valuation cube selector and records the new logarithmic-subfield candidate.

\section{Confirmation of current spine}
The affine base, residual contraction, cotree witness bijection, and upper-bound escapes remain PROVED/EXACT. The Gold-cube zero-kernel firewall (CUBEPERMFW1), the planted even-cube section firewall (PLANTEDCUBEFW1), CUBEROOTSEL1 ($\nu_3(d)=0$), and the fixed-$\nu_3(d)$ subfield selector (VALSUBCUBEFW1) are retained as PROVED/EXACT after independent reviews. The logarithmic-subfield selector for even degrees satisfying $3^{\nu_3(d)}=O(\log d)$ has a complete proof in the synthesis snapshot and is retained as CANDIDATE pending independent specialist review. Growing-factor amplification remains CONDITIONAL. The unrestricted whole-computation lower bound remains OPEN.

VALSUBCUBEFW1 and its synthesis promotion correctly establish that for each fixed $a=\nu_3(d)\ge1$, the subfield $\mathbb F_{2^{2\cdot3^a}}\subseteq\mathbb F_{2^d}$ supplies a cubic nonresidue via constant-size enumeration (for fixed $a$), enabling an Adleman--Manders--Miller correction that selects one cube root from the target alone in polynomial time. The square-cotree case $K_{7,7}$ with $d=36$ is covered. The 7,659-assertion audit is COMPUTATIONAL and matches. Thus fixed valuation, including genuine square cotree geometry, is FALSIFIED/RETIRED as a sufficient residual-work premise.

The logarithmic extension enlarges the firewall to slowly unbounded valuations under the stated bound; its proof is elementary and correct on the displayed arithmetic, but remains CANDIDATE as newly introduced in this cycle. The complementary regime where $3^{\nu_3(d)}/\log d$ is unbounded is the next algorithmic target; failure of known selectors is not a lower bound.

\section{Open load-bearing arrows}
The implication from an arbitrary eligible short program (or uniform decider) to superpolynomial cost on a structured residual object remains OPEN. No model-matched all-algorithm theorem exists. The firewalls relativize and algebrize; no nonrelativizing hinge is known. Component ranks, polar spectra, fiber cardinalities, collision-effective entropy, graph treewidth, and the mere absence of a total inverse or fixed-valuation selector are not runtime theorems. The next necessary test is selector hardness on the super-logarithmic-subfield degrees, not merely lack of one particular algorithm.

\section{Highest-value next target}
The single highest-value next assignment stated in the living document remains appropriate: independently review the logarithmic-subfield theorem, then attack the complementary degrees where $3^{\nu_3(d)}/\log d$ is unbounded using Sze-style nonresidue-free extraction and explicit $3$-power subfield towers. A successful selector is a useful firewall; absence of a recognized normal form is not hardness.

\section{Status and recommendations}
P versus NP remains open. The living document is updated to record the independent confirmation of the fixed-valuation promotion and the retention of the logarithmic-subfield theorem as CANDIDATE. All claims continue to be classified under the allowed statuses. The comments folder receives this review artifact.

% END_OF_VERIFIED_REVIEW::GROK::GROK-0114
\end{document}
